\documentclass[11pt,a4paper]{article}
\usepackage{fontspec,amsmath,amssymb,amsthm,microtype}
\setmainfont{lmroman10-regular.otf}[BoldFont=lmroman10-bold.otf,ItalicFont=lmroman10-italic.otf,BoldItalicFont=lmroman10-bolditalic.otf]
\setmonofont{lmmono10-regular.otf}[Scale=MatchLowercase]
\usepackage[margin=25mm]{geometry}
\usepackage[unicode,hidelinks]{hyperref}
\setlength{\emergencystretch}{3em}
\setlength{\parskip}{0.4em}
\setlength{\parindent}{0pt}
\allowdisplaybreaks[2]
\raggedbottom
\title{\texorpdfstring{Proximal Langevin library entry point}{Proximal Langevin library entry point}}
\author{LangevinLean mathematical exposition}
\date{}
\begin{document}\maketitle
\noindent Source module: \texttt{Proximal\allowbreak{}Langevin.\allowbreak{}lean}. The accompanying HTML pairs this mathematical exposition with the exact Lean statements and records the scope of the formalization.


The root module gathers proximal calculus, actual particle algorithms, corrected original-center bounds and constructed mean-field comparisons into one importable Lean library. Its imported theorems retain their exact assumptions and scope.


\section{\texorpdfstring{Follow the imported proof families and reproduce the build}{Follow the imported proof families and reproduce the build}}


Importing \texttt{Proximal\allowbreak{}Langevin} makes the checked public development available. The mathematical chain is true proximal minimization, residual/\allowbreak{}Moreau calculus, exact algorithm identities, genuine finite-moment Gaussian laws, concentration and transport, then the original-center comparison. The splitting chain additionally constructs Brownian stationarity and orthogonality; the mean-field branch constructs continuous \(L^2\) integral equations and actual iid comparisons. Run \texttt{lake build Proximal\allowbreak{}Langevin +Proximal\allowbreak{}Langevin.\allowbreak{}Check}. This umbrella has no local declarations; its proof content is in the linked imported modules.



This entry module assembles the named proof families through imports. The complete assumptions, intermediate estimates and final conclusions are stated in the dependency pages below. Building the entry module recompiles its imported dependencies and verifies that their interfaces agree.


\end{document}
